For the sake of completeness we give here a full description of the NISC-based PPB originally described in \cite{EC:KohLysNgu23} when instantiated using oblivious transfer and garbled circuits.
As the proofs of correctness and security are provided in \cite{EC:KohLysNgu23} we omit them here.

Before describing the instantiation we first briefly recall the interfaces of oblivious transfer and garbling schemes.

\begin{definition}[Oblivious Transfer]
	A one-out-of-two Oblivious Transfer (OT) protocol is a tuple of algorithms $\ot = (\otinit, \otsend, \otreceive)$ defined as:
	\begin{itemize}
		\item $\otinit(\secparam, b)$ takes as input a bit $b \in \bin$ and outputs a secret key/public key pair $(\otseck, \otpubk)$.
		\item $\otsend(\otpubk, \otval_0, \otval_1)$ takes as input a public key and two values and outputs a ciphertext $\otct$.
		\item $\otreceive(\otseck, \otct)$ takes as input a secret key and a ciphertext and outputs a value $\otval'$.
	\end{itemize}

	We refer the reader to other texts discussing oblivious transfer in more depth such as \cite{EC:Kalai05} for a complete definition.
\end{definition}


\begin{definition}[Garbling Schemes]
	A garbling scheme is a tuple of algorithms $\gc = (\gcircuit,$ $\ginput, \geval)$ defined as:
	\begin{itemize}
		\item $\gcircuit(\secparam, \circuit)$ takes as input a circuit $\circuit: \bin^n \rightarrow \bin^m$ and outputs the garbled circuit $\gcirc$ and key $k$.
		\item $\ginput(k, x)$ takes as input a key $k$ and string $x \in \bin^n$ and outputs a garbled input $\tilde{x}$.
		\item $\geval(\gcirc, \tilde{x})$ takes as input a garbled circuit and input, and outputs $y \in \bin^m$.
	\end{itemize}

	We say that a garbling scheme is \emph{projective} if each bit of the garbled input $\gx$ depends on one bit of the actual input $x$, i.e. each bit of the input is garbled independently of the other bits.
	We will let $\pginput(k, x_i)$ denote the algorithm to garble a single bit of input, i.e.
	$\ginput(k, x) = \{\pginput(k, i, \xthbit{x}{i})\}_{i \in [n]}$ for $x \in \bin^n$.

	As with OT, we refer the reader to \cite{CCS:BelHoaRog12} for a full description of garbling schemes.
\end{definition}

\subsection{Construction}

Define the following two relations:
\begin{align*}
	\kgrel^{\lambda, \comparams, \numdigits} = & \left\{
	((pk', \com), (\thr, sk', r_k, r_c)) \::\:
	\begin{gathered}
		(sk', pk') = \kgen'(\lambda, \thr; r_k) \\
		\wedge \: \com = \commit(\comparams, \thr; r_c)
	\end{gathered}
	\right\}                                             \\
	\erel^{\lambda, \comparams, \numdigits} =  & \left\{
	((Z', \com, pk'), (\uval, \msg, r_e, r_c)) \::\:
	\begin{gathered}
		Z' = \escrow'(\secpar, pk', (\uval, \msg); r_e) \\
		\wedge \: \com = \commit(\comparams, (\uval, \msg); r_c)
	\end{gathered}
	\right\}
\end{align*}

Let $\kglang^{\lambda, \comparams}$ and $\elang^{\lambda, \comparams}$ be the corresponding languages to these relations, and
$\kgps^{\lambda, \comparams}$, $\eps^{\lambda, \comparams}$ be corresponding NIZKs for these languages.
We may drop the superscripts when the values are clear from context.
The generic NISC-based PPB construction instantiated with garbled circuits and OT can be seen in \cref{fig:PPB_generic_construction} and \cref{fig:PPB_generic_construction_cont}.

\begin{figure}
	\begin{pchstack}[boxed, center, space=1em]

		\begin{pcvstack}[space=1em]
			\begin{pchstack}[center, space=1em]
				\procedure[linenumbering]{$\setup(\secparam, \comparams)$}{
					(\crs_{KG}, \td_{KG}) \leftarrow \zksetup_{KG}(\secparam) \\
					(\crs_{E}, \td_{E}) \leftarrow \zksetup_{E}(\secparam) \\
					\ppbpp = (\secpar, \comparams, \crs_{KG}, \crs_{E}) \\
					\pcreturn(\ppbpp, (\td_{KG}, \td_{E}))
				}

				\procedure[linenumbering]{$\kgen'\left(\lambda, x\right)$}{
					n := |x| \\
					\pcfor i \in [n] \pcdo \\
					\t (\otseck_i, \otpubk_i) \leftarrow \otinit(\secparam, x[i])\\
					sk' \leftarrow \{\otseck_i\}_{i \in [n]} \\
					pk' \leftarrow \{\otpubk_i\}_{i \in [n]} \\
					\pcreturn (sk', pk')
				}
			\end{pchstack}

			\procedure[linenumbering]{$\kgen\left(\ppbpp, x, r_{KG}\right)$} {
				\pcparse (\secpar, \comparams, \crs_{KG}, \crs_{E}) = \ppbpp \\
				(sk', pk') \xleftarrow{r_k} \kgen'(\lambda, x) \\
				\com_{KG} \leftarrow \commit(\comparams, x; r_{KG}) \\
				\zkproof_{KG} \leftarrow \zkprove_{KG}(\crs_{KG}, (pk', C), (\thr, sk', r_k, r_c)) \\
				pk \leftarrow (pk', \com_{KG}, \zkproof_{KG}) \\
				sk \leftarrow (sk', pk) \\
				\pcreturn (sk, pk)
			}

			\procedure[linenumbering]{$\verpk(\ppbpp, pk, \com_{KG})$}{
				\pcparse (\secpar, \comparams, \crs_{KG}, \crs_{E}) = \ppbpp \\
				\pcparse (pk', \com'_{KG}, \zkproof_{KG})  = pk \\
				\pcreturn \zkverify_{KG}(\crs_{KG}, \zkproof_{KG}, (pk', \com_{KG})) \\
				\t \wedge (\com'_{KG} = \com_{KG})
			}

			\procedure[linenumbering]{$\verescrow(\ppbpp, pk, \com_{E}, Z = (Z', \zkproof_{E}))$}{
				\pcparse (\secpar, \comparams, \crs_{KG}, \crs_{E}) = \ppbpp \\
				\pcparse (pk', \com'_{KG}, \zkproof_{KG})  = pk \\
				\pcreturn \verpk(\ppbpp, pk, \com_{KG}) \\
				\t \zkverify_E((Z', \com_E, pk'), \zkproof_E)
			}
		\end{pcvstack}


	\end{pchstack}
	\caption{A generic NISC-based PPB construction}
	\label{fig:PPB_generic_construction}
\end{figure}

\begin{figure}
	\begin{pcvstack}[boxed, center, space=1em]
		\procedure[linenumbering]{$\escrow'(\secpar, pk', \uval)$}{
			\pcparse pk' = \{\otpubk_i\}_{i \in [n]} \\
			(k, \gcirc) \leftarrow \gcircuit(\secparam, C) \\
			\pcfor i \in [n] \pcdo \\
			\t \otct_i \leftarrow \otsend(pk_i, \pginput(k, i, 0), \pginput(k, i, 1)) \\
			\pcfor i \in [|\uval|] \pcdo \\
			\t \gx_{n + i} \leftarrow \pginput(k, n + i, \xthbit{\uval}{i}) \\
			ct' := \{ct_1,\dots, ct_n\} \\
			\gx' := \{\gx_{n + 1}, \dots, \gx_{n + |\uval|}\} \\
			\pcreturn (\gcirc, ct', \gx')
		}

		\procedure[linenumbering]{$\escrow(\ppbpp, pk, \uval, r_E)$}{
			\pcparse (\secpar, \comparams, \crs_{KG}, \crs_{E}) = \ppbpp \\
			\pcparse (pk', \com_{KG}, \zkproof_{KG}) = pk \\
			Z' \xleftarrow{r_e} \escrow'(\secpar, pk', \uval) \\
			\com_{E} \leftarrow \commit(\comparams, \uval; r_E) \\
			\zkproof_E \leftarrow \zkprove_E(\crs_E, (Z', \com_E, pk'), (\uval)) \\
			\pcreturn (Z', \zkproof_E)
		}

		\procedure[linenumbering]{$\dec(\ppbpp, sk, \com_E, Z = (Z', \zkproof_E))$}{
			\pcparse (\secpar, \comparams, \crs_{KG}, \crs_{E}) = \ppbpp \\
			\pcparse (\{\otseck_i\}_{i \in [n]} , pk) = sk \\
			\pcif \verescrow(\ppbpp, pk, \com_E, Z) = 0 \pcthen \pcreturn \bot \\
			\pcparse (\gcirc, \{\otct_1,\dots,\otct_n\}, \{\gx_{n+1},\dots,\gx_{n + |\uval|}\}) = Z' \\
			\pcfor i \in [n] \pcdo \\
			\t \gx_i \leftarrow \otreceive(\otseck_i, \otct_i) \\
			\gx := \{\gx_i\}_{i \in [n + |\uval|]} \\
			\pcreturn \geval(\gcirc, \gx)
		}
	\end{pcvstack}
	\caption{A generic NISC-based PPB construction, continued.}
	\label{fig:PPB_generic_construction_cont}
\end{figure}



